\chapter{Cơ sở lý thuyết và công trình liên quan}
\label{chap:related-work}

\section{Suy luận tại thời điểm kiểm thử}

Trong khóa luận này, \emph{test-time compute} là lượng tính toán được sử dụng sau khi
mô hình đã được huấn luyện, trong quá trình sinh lời giải cho một đầu vào cụ thể. Cách
tiếp cận này không thay đổi trọng số của mô hình nền. Thay vào đó, hệ thống điều chỉnh
số token tối đa cho một quỹ đạo sinh, số quỹ đạo được lấy mẫu, hoặc thời điểm dừng lấy
mẫu. Chain-of-thought cho phép mô hình biểu diễn các bước trung gian trước đáp án
cuối~\cite{wei2022cot}; đây là tiền đề để tăng compute suy luận có thể làm thay đổi xác
suất giải đúng một bài toán.

Khóa luận phân biệt hai trục điều khiển. \emph{Breadth} là số quỹ đạo độc lập $k$.
\emph{Depth} là cách gọi tắt cho giới hạn sinh $L=\code{max\_new\_tokens}$ của mỗi
quỹ đạo. Depth trong tài liệu này không hàm ý rằng lời giải dài hơn luôn có cấu trúc
suy luận sâu hơn. Hai trục có tác động khác nhau: tăng $L$ tạo cơ hội để một quỹ đạo
hoàn tất, còn tăng $k$ tạo thêm các lần thử và cho phép tổng hợp đáp án. Vì chi phí
thực tế của một quỹ đạo có thể thấp hơn giới hạn, đánh giá phải dùng số token thật sự
đã sinh thay vì chỉ dùng tích danh nghĩa $kL$.

\section{Lấy nhiều mẫu và self-consistency}

Self-consistency lấy nhiều quỹ đạo chain-of-thought rồi chọn đáp án xuất hiện thường
xuyên nhất~\cite{wang2023selfconsistency}. Repeated sampling mở rộng trực giác này:
khi các lần sinh có tính ngẫu nhiên, nhiều mẫu làm tăng xác suất ít nhất một mẫu tìm
được lời giải đúng~\cite{brown2024monkeys}. Tuy nhiên, tăng số mẫu đồng đều cho mọi
bài làm chi phí tăng tuyến tính và không tận dụng việc một số bài đã có thể dừng sớm.
Ngoài ra, nếu mọi mẫu đều bị cắt trước đáp án cuối vì $L$ quá thấp thì tăng breadth
không giải quyết được nút thắt token cap.

Khóa luận sử dụng biểu quyết chỉ trên các đáp án cuối parse được. Mẫu không hợp lệ
không được bỏ phiếu và trường hợp không có đáp án hợp lệ được tính là sai. Quy tắc này
khắt khe hơn việc trích số cuối cùng xuất hiện trong toàn bộ reasoning, nhưng tránh biến
một con số trung gian thành dự đoán. Khi hai đáp án có cùng số phiếu, đáp án thuộc mẫu
có chỉ số thấp hơn được chọn để bảo đảm tính tất định.

\section{Thích ứng theo breadth và dừng sớm}

Một hướng nghiên cứu giảm chi phí self-consistency là quyết định số mẫu theo tín hiệu
đã quan sát. Adaptive-Consistency ước lượng mức ổn định của phân phối đáp án để dừng
lấy mẫu~\cite{aggarwal2023adaptive}. Difficulty-Adaptive Self-Consistency điều chỉnh
chi phí theo độ khó dự đoán của đầu vào~\cite{wang2025penny}, còn Reasoning-Aware
Self-Consistency sử dụng thông tin từ các đường suy luận để lấy mẫu hiệu quả
hơn~\cite{wan2025rasc}. ESC nghiên cứu dừng self-consistency sớm~\cite{li2024esc}, và
Certaindex xây dựng chỉ số độ chắc chắn để mở rộng suy luận hiệu quả
hơn~\cite{fu2025certaindex}.

Các phương pháp trên cho thấy không nhất thiết phải trả chi phí của một $k$ cố định
cho mọi bài. Tuy nhiên, một số phương pháp cần ước lượng độ chắc chắn, đặc trưng phức
tạp hoặc nhiều mẫu hơn phạm vi phần cứng của đề tài. Baseline breadth-only của khóa
luận vì vậy dùng một quy tắc tối giản và có thể kiểm toán: luôn lấy hai mẫu tại một $L$
cố định; dừng nếu cả hai đáp án hợp lệ và giống nhau; nếu không thì lấy đủ bốn mẫu.
Quy tắc này dùng đồng thuận đáp án chuẩn hóa chứ không giả định đồng thuận đồng nghĩa
với đúng. Nó là một chính sách chi phí, không phải một bộ ước lượng xác suất đúng.

\section{Thích ứng theo ngân sách token}

Nhánh thứ hai điều khiển độ dài suy luận. Token-Budget-Aware reasoning nghiên cứu việc
đưa ràng buộc ngân sách vào quá trình suy luận~\cite{han2025token}. Thought Calibration
điều chỉnh compute kiểm thử dựa trên độ tin cậy~\cite{wu2025thought}. Các công trình
về compute-optimal inference cho thấy điểm vận hành tối ưu phụ thuộc vào bài toán,
mô hình và loại tài nguyên được cấp~\cite{snell2025scaling,wu2024inference}. Learning
How Hard to Think xem phân bổ compute theo đầu vào như một bài toán thích ứng
riêng~\cite{damani2025hard}.

Một số hướng can thiệp trực tiếp vào quá trình sinh. s1 đề xuất budget forcing để kéo
dài hoặc ép quá trình suy luận theo ngân sách~\cite{muennighoff2025s1}; L1 học cách
điều khiển độ dài suy luận bằng reinforcement learning~\cite{aggarwal2025l1}. Ngược
lại, ``When More is Less'' cảnh báo rằng chain-of-thought dài hơn không mặc nhiên tạo
ra kết quả tốt hơn~\cite{wu2025moreless}. Các kết quả này là cơ sở để khóa luận không
đồng nhất token nhiều với reasoning tốt và luôn báo cáo cả accuracy lẫn token thực tế.

Phạm vi thực nghiệm chính sử dụng decoder Qwen3 không chỉnh sửa với thinking mode
được bật. Non-thinking mode được đánh giá bổ sung trên Validation sau khi Test chính
đã đóng băng. Maximum-budget force-answer được đánh giá bằng protocol bổ sung
sau Validation tại các hard cap; native numeric thinking budget không có trong
backend cố định. Phân tách này giữ so sánh policy chính ở cùng cơ chế sinh và
không dùng kết quả decoding-control để chọn lại policy.

\section{Phân bổ kết hợp và chính sách học được}

Phân bổ joint phải trả lời hai câu hỏi: chọn token cap nào và có nên lấy thêm mẫu hay
không. Các phương pháp dựa trên độ khó hoặc tối ưu chính sách có thể học quyết định
compute từ dữ liệu~\cite{wang2025penny,zhai2026adaptive}. PaCoRe mở rộng theo hướng
phối hợp reasoning song song và tuần tự~\cite{hu2026pacore}. Những hướng này giàu khả
năng hơn một heuristic, nhưng đòi hỏi dữ liệu huấn luyện controller đủ lớn, định nghĩa
reward rõ ràng và một quy trình chống rò rỉ nghiêm ngặt.

Trong dữ liệu Train của đề tài, số sự kiện escalation dương tại từng stage không đạt
ngưỡng tối thiểu đã đăng ký là 30. Vì vậy, một joint controller học máy đáng tin cậy
chưa được xem là đóng góp đã hoàn tất. Chính sách joint được đánh giá trong khóa luận
là heuristic tất định: prompt-token tertile từ Train chọn $L$; sau đó policy dừng ở
$k=1$ chỉ khi đáp án cuối hợp lệ và mẫu không bị cắt, dừng ở $k=2$ khi hai đáp án hợp
lệ giống nhau, nếu không lấy đủ $k=4$. Việc gọi đúng tên heuristic là cần thiết để
không phóng đại mức độ mới của phương pháp.

\section{Hiệu quả suy luận cho mô hình nhỏ}

Reasoning in Token Economies nhấn mạnh rằng đánh giá chiến lược suy luận phải gắn với
ngân sách thay vì chỉ so sánh độ chính xác~\cite{wang2024tokeneconomies}. Qwen3 cung
cấp thinking mode và nhiều kích thước mô hình, tạo điều kiện nghiên cứu test-time
scaling trên mô hình mở~\cite{qwen3report}. Khóa luận chọn Qwen3-0.6B vì có thể chạy
lặp lại trên một RTX 4060. Phân tích Qwen3-1.7B trước đó được giữ như bằng chứng lịch
sử nhưng không gộp với kết quả chính, do khác kích thước mô hình và action space.

Việc dùng mô hình nhỏ làm cho bài toán phân bổ compute trở nên thực tế: một cấu hình
quá lớn có thể tăng thời gian chạy đáng kể mà không cải thiện accuracy tương ứng.
Đồng thời, kết luận không được khái quát sang mọi họ mô hình. Trong phạm vi khóa luận,
``hiệu quả'' được hiểu là vị trí trên quan hệ accuracy--actual-generated-token của
cùng mô hình, cùng dữ liệu và cùng quy trình sinh.

\section{Khoảng trống nghiên cứu và vị trí của đề tài}

Các công trình trước đã nghiên cứu self-consistency, dừng sớm, token budget và phân bổ
compute học được. Khoảng trống mà đề tài xử lý ở quy mô khóa luận đại học không phải
là đề xuất một họ thuật toán hoàn toàn mới, mà là đánh giá kiểm soát một câu hỏi hẹp:
trên mô hình mở 0.6B chạy cục bộ, một chính sách tuần tự không verifier kết hợp token
cap và selective sampling có cải thiện biên Pareto so với các đối chứng cố định,
depth-only và breadth-only hay không.

Vị trí thực nghiệm của đề tài được xác định bởi năm ràng buộc. Thứ nhất, hệ thống không
dùng verifier ngoài hoặc hidden state. Thứ hai, đáp án cuối được parse nghiêm ngặt và
mẫu không hợp lệ không được bỏ phiếu. Thứ ba, chi phí là token thật sự sinh và mọi mẫu
đã yêu cầu đều bị tính phí. Thứ tư, policy được chọn trên Train/Validation rồi đóng
băng trước khi mở Test. Thứ năm, oracle chỉ là cận trên cơ hội từ sample pool đã sinh,
không phải policy triển khai được. Các ràng buộc này tạo một phép so sánh nhỏ nhưng
có thể tái lập, đồng thời làm rõ giới hạn của mọi kết luận rút ra.

GSM8K~\cite{cobbe2021gsm8k} và MATH-500, được lấy từ bộ MATH
\cite{hendrycks2021math}, cung cấp hai mức độ và kiểu đáp án khác nhau. Việc một policy
hoạt động trên GSM8K nhưng không ổn định trên MATH-500 được xem là bằng chứng về tính
phụ thuộc miền, thay vì được che khuất bằng cách gộp hai tập dữ liệu. Chương tiếp theo
chuyển các nguyên tắc trên thành định nghĩa chính sách, cost accounting và protocol
thống kê cụ thể.
